\section{Reset: an intrinsic binary fragment}\label{sec:reset}

The derived operation
\begin{equation}\label{reset:eq:reset}
 A\del B=\begin{cases}B,&|B|=1,\\A\cup B,&|B|\geq2\end{cases}
\end{equation}
retains a resolved right operand and otherwise unions. It is realized by
starting with $J(A,B)$ and applying $U\mapsto U\star_r B$ once for each
role. If $B=\{r\}$, that owner's suppression selects it and all later
owners preserve it. If $B$ is ambiguous, every step preserves $A\cup B$.
The order of the suppressions is immaterial.

Reset commutes with every role permutation. Its language consists only
of variables and $\del$. It cannot express arbitrary owned choice or
ordinary union: reset preserves singleton-valued inputs, whereas an
owned choice between two other winners can return their pair.
We first show why this particular fragment is intrinsic, then describe
its terms, known equational basis, and exact prerequisite semantics.

\subsection{An intrinsic binary fragment}

The operation also has a characterization inside owned choice.
A term \emph{preserves atoms} if it returns a singleton whenever all
inputs are singletons. Work first with constant-free, switch-free
owned-choice terms, which are automatically idempotent.

\begin{lemma}[One-constant rigidity]\label{reset:lem:rigidity}
A unary owned-choice expression whose only constant is $\{r\}$ is
either constantly $\{r\}$ or inflationary: it contains its input on
every valuation. Consequently, returning $\{r\}$ at any other singleton
$\{q\}$ already forces it to be constant.
\end{lemma}
\begin{proof}
Induct on expressions. Combining two inflationary expressions preserves
inclusion of the input: if an owner wins outright, that singleton
already contains the input. Combining an inflationary expression with
the constant $\{r\}$ returns the constant when the owner is $r$;
otherwise it remains inflationary, because the constant operand cannot
give a different owner an outright win. Two constant expressions remain
constant. An inflationary function cannot send $\{q\}$ to $\{r\}$ when
$q\ne r$. The dichotomy and this rigidity consequence are
kernel-checked in \texttt{OneConstant.lean}.
\end{proof}

\begin{theorem}[Exact binary atom-preserving fragment]\label{reset:thm:binary}
For $K\geq3$, the binary atom-preserving operations definable by
constant-free owned choices are exactly
\[
 x,\qquad y,\qquad x\del y,\qquad y\del x.
\]
Hence the reset clone is exactly the part generated by the binary
operations of the idempotent, atom-preserving owned-choice clone.
\end{theorem}
\begin{proof}
The atomic restriction of an idempotent atom-preserving owned term is
a projection. This is \cref{gen:thm:S2}, using $\operatorname{Pol}(\iota_3)=S_2$
\cite[Proposition~2.2]{Szendrei2012}, and an idempotent operation on
$K\geq3$ atoms has image size at least three, so the essentially unary
case of $S_2$ must apply.

Suppose the atomic restriction projects $x$. If $y$ does not occur,
idempotency makes the whole term $x$. Otherwise both variables occur.
Freeze $x=\{r\}$, leaving a unary expression in $y$ with only that
constant. Its atomic outputs are $\{r\}$, so
Lemma~\ref{reset:lem:rigidity} makes it constant on all $y$.

Now let $x$ be ambiguous. If $y$ is ambiguous too, every node computes
union, giving $x\cup y$. If $y=\{r\}$, specialize instead in $x$.
Its atomic restriction is the identity, so rigidity makes it
inflationary. Its result is contained in $x\cup\{r\}$, since owned
choice never introduces a winner absent from its inputs. It contains
$r$ as well: under this valuation every singleton subexpression is
$\{r\}$, by induction, so no different owner can suppress the
contribution of an occurring $y$. The value is therefore
$x\cup\{r\}$. These cases give exactly $y\del x$.
The opposite projection gives the other two operations, all of which
are definable.
\end{proof}

The idempotent-elimination theorem (\cref{gen:thm:core})
extends this characterization to idempotent full-signature terms,
including constants and internal switches in their original syntax.
The binary theorem uses that mathematical theorem and the
external atomic-clone classification; its complete proof is not yet a
Lean theorem.

The binary restriction is strict. For an owner $a$, the ternary operation
\[
 H_a(x,y,z)=(y\star_a z)\del x
\]
is idempotent and atom-preserving. Set $x=\{a,b\}$ and $z=\{c\}$.
At $y=\{a\}$ its value omits $c$, while at $y=\{b\}$ it contains $c$.
The ambiguity statuses are unchanged, and only $z$ supplies $c$.
The selector semantics below therefore rules out a reset expression
for $H_a$. Thus reset is an intrinsic binary core, while higher-arity
atom-preserving operations still retain owner-sensitive behavior.

\subsection{Exact prerequisite semantics}

Let $N$ be the finite set of input positions, including any declared
unused positions. For a valuation $S=(S_j)_{j\in N}$ put
$B=\{j:|S_j|\geq2\}$. This is a set of input positions, not winner roles.
The root $\rho(t)$ is its rightmost variable:
\[
 \rho(x_j)=j,\qquad \rho(s\del t)=\rho(t).
\]

\begin{lemma}[Resolution and selection]\label{reset:lem:selection}
A reset term is ambiguous exactly when its root input is ambiguous.
There is a selector $E_t:2^N\to2^N$ with
\begin{align*}
 E_{x_j}(B)&=\{j\},\\
 E_{s\del t}(B)&=E_t(B)\cup
 \begin{cases}E_s(B),&\rho(t)\in B,\\\varnothing,&\rho(t)\notin B,\end{cases}\\
 t(S)&=\bigcup_{j\in E_t(B)}S_j.
\end{align*}
Every selector is monotone, always selects its root, and selecting
position $j$ is independent of whether $j$ itself belongs to $B$.
\end{lemma}
\begin{proof}
If the right child is resolved, reset returns it. If it is ambiguous,
reset unions the children and remains ambiguous. This proves the first
three statements by induction. Monotonicity follows from the positive
recursion. For self-independence, at a node with root $j$ that input is
always selected. Otherwise its new guard concerns a different position,
and both child selection predicates are independent of $j$ by induction.
\end{proof}

For each position $j$, record the inclusion-minimal sets $A\subseteq N$
such that $j\in E_t(A)$. Denote this antichain by $\mathcal G_j(t)$.
An empty antichain means the position is unused. Then
\begin{equation}\label{reset:eq:guards}
 j\in E_t(B)\quad\Longleftrightarrow\quad
       \exists A\in\mathcal G_j(t)\; A\subseteq B.
\end{equation}
For the root its antichain is $\{\varnothing\}$, and every other
antichain omits $\varnothing$. No minimal prerequisite for $j$ contains
$j$, by self-independence.

These tables are computed directly. A variable contributes its empty
guard. At $s\del t$, retain the right table and add $\rho(t)$ to each
left guard, then delete duplicates and every guard containing a smaller
one for the same output. Before minimization there is at most one guard
per variable occurrence. A tree traversal can collect these guards and
normalize them in polynomial time.

\begin{theorem}[Canonical equality and three-element completeness]
\label{reset:thm:three}
Two reset terms are equal on all nonempty-set valuations if and only if
their prerequisite tables agree. The result is independent of $I$ for
every $|I|\geq3$. It suffices to evaluate on a fixed three-element chain
\[
 Z=\{a\}\subset U=\{a,b\}\subset T=\{a,b,c\}.
\]
Consequently equality is decidable in polynomial time in reset-tree size.
A failed equality has a polynomially constructible countervaluation
using only $Z,U,T$.
\end{theorem}
\begin{proof}
Equal tables give equal selectors and therefore equal values.
If two tables differ, some selector inclusion differs at $B,j$.
Assign $T$ to position $j$, $U$ to the other positions in $B$, and $Z$
elsewhere. Setting $j$ to ambiguous does not change whether $j$ is
selected, by Lemma~\ref{reset:lem:selection}. The marker $c$ occurs in the
result exactly when that position is selected. Thus the terms differ.
Distinct minimal antichains give such $B,j$: take a minimal prerequisite
on one side not containing a prerequisite from the other. All operations
on finite guard tables are polynomial.
\end{proof}

The chain is closed under reset. Numbering it $0<1<2$, its table is
\begin{equation}\label{reset:eq:three-table}
\begin{array}{c|ccc}
\del&0&1&2\\\hline
0&0&1&2\\
1&0&1&2\\
2&0&2&2
\end{array}
\end{equation}
This is the dual of $G_3$ in Waldhauser's classification
\cite[Theorem~3.4, table $G_3$]{Waldhauser2006}, under
$0\mapsto b$, $1\mapsto a$, $2\mapsto c$.
This identifies the reset variety with the variety generated by that
established three-element algebra.
It is also exactly the order algebra $O_{2+1}$ of
Freese et al.\ \cite[Section~3]{FreeseEtAl2002}. Their poset has $a<b$
and an isolated $c$, with $x\cdot y=x$ if $x\leq y$ and $y$ otherwise.
The isomorphism from our labels is $0\mapsto c$, $1\mapsto b$,
$2\mapsto a$. This identification gives a finite basis immediately.

\begin{corollary}[Minimum finite generator]
The three-element algebra generates the common variety of all reset
algebras on $\Pow(I)$ for $|I|\geq3$, and no smaller algebra generates
that variety.
\end{corollary}
\begin{proof}
Theorem~\ref{reset:thm:three} identifies their identities. The three-element
algebra is nonassociative:
$(2\del0)\del1=1$, while $2\del(0\del1)=2$.
Every idempotent binary operation on two elements is one of the two
projections or two semilattice operations; all four are associative.
Any generator would have to satisfy the valid identity $x\del x=x$,
so a smaller generator is impossible.
\end{proof}

This is a bound on the carrier size of a generator, not on the arities
available in a clone on a larger carrier. The absence of intermediate
nontrivial subclones of the reset clone follows instead from the
identification-minor argument in \cref{reset:thm:minimal}. It does not
classify clones between reset and the larger owned-choice clone.

Even a two-value input domain that is not closed under reset cannot
test all identities: associativity holds whenever its three arguments
use at most two distinct winner sets. A failure requires the middle
argument to be atomic and the last ambiguous, with a contribution of
the first lying outside their union. Two distinct input values cannot
supply that contribution.

\subsection{A finite equational calculus from order algebras}

\begin{theorem}[Known reset basis]\label{reset:thm:basis}
For every $K\geq3$, an equation between reset terms holds on all
nonempty winner-set valuations if and only if it follows in equational
logic from
\begin{align*}
 x\del x&=x,\tag{R1}\\
 (x\del y)\del x&=y\del x,\tag{R2}\\
 (x\del y)\del y&=x\del y,\tag{R3}\\
 x\del((x\del y)\del z)&=x\del(y\del z),\tag{R4}\\
 ((x\del y)\del z)\del y&=(x\del z)\del y,\tag{R5}\\
 x\del(y\del z)&=y\del(x\del z).\tag{R6}
\end{align*}
\end{theorem}
\begin{proof}
Freese et al.\ \cite[Theorem~3.7]{FreeseEtAl2002} prove that their
equations (E1)--(E5), together with left commutation, form a basis for
$O_{2+1}$. These are exactly (R1)--(R6) under the displayed isomorphism.
Theorem~\ref{reset:thm:three} identifies the identities of that algebra with
those of winner-set reset for every $K\geq3$.
\end{proof}

The inference rules are the ordinary rules of equational logic:
reflexivity, symmetry, transitivity, substitution, and congruence under
$\del$. This is an application of the published basis, not a new basis.
Soundness of its six equations is kernel-checked in \texttt{Reset.lean};
the completeness argument uses the external theorem and has not been
mechanized. This result concerns reset, not the larger owned-choice
signature.

The same source already develops common-root path normal forms and a
suffix criterion for comparing them
\cite[Lemmas~3.5--3.6]{FreeseEtAl2002}. The correspondence with guards
is exact. For a left-associated path $p$ with pairwise distinct input
positions $i_1,\ldots,i_m$, let
\[
 \sigma_p(i_a)=\{i_{a+1},\ldots,i_m\}.
\]
Its root is $i_m$. Unfolding the selector recursion gives
$i_a\in E_p(B)$ iff $\sigma_p(i_a)\subseteq B$: an occurrence
contributes precisely when every variable to its right is ambiguous.
If $w$ is a product of paths with that common root, its selectors are
the unions of the path selectors. Its canonical guards are therefore
\[
 \mathcal G_j(w)=\min_{\subseteq}
     \{\sigma_p(j):p\text{ is a constituent path of }w,\ j\in p\}.
\]
For another path $p$ with the same root, requiring
$E_p(B)\subseteq E_w(B)$ for every $B$ is equivalent to requiring,
for each occurrence of $j$ in $p$,
some suffix for $j$ in $w$ contained in $\sigma_p(j)$. This is the
suffix-inclusion criterion of their Lemma~3.6, expressed through our
monotone selection predicates. Their order inequality $w\leq p$ reads
$w\del p=w$; for common-root terms this is exactly the displayed
inclusion of selectors. Thus the guard tables retain the minimal
suffixes needed for evaluation and comparison; they do not introduce
an unrelated normalization theory. The following sections give the
prerequisite/antimatroid presentation and its algorithms. Priority for
those connections and algorithms remains unestablished.

\subsection{The antimatroid represented by a term}

Prerequisites can describe alternative routes. For example, let
\[
 t=((x_3\del x_1)\del x_0)\del((x_3\del x_2)\del x_0).
\]
Its canonical table is
\[
 \mathcal G_0=\{\varnothing\},\quad
 \mathcal G_1=\mathcal G_2=\{\{0\}\},\quad
 \mathcal G_3=\{\{0,1\},\{0,2\}\}.
\]
Input $0$ is always retained. Inputs $1,2$ are retained when $0$ is
ambiguous, while input $3$ additionally requires either $1$ or $2$ to
be ambiguous. This alternative cannot be described by a single
partial order of prerequisites. Feasible sets record which positions
can be reached by following the alternatives.

We use the standard finite feasible-set definition
\cite[Section~7]{Czedli2014}. An antimatroid on its actual support $X$
is a union-closed family $\mathcal F\subseteq2^X$ containing
$\varnothing$ and $X$, with the accessibility property: every nonempty
feasible set can lose one element and remain feasible. We require
exactly one feasible singleton $\{r\}$. Accessibility then makes $r$
belong to every nonempty feasible set. Unused input positions lie
outside $X$.

Write
\[
 \bas_{\mathcal F}(B)=\bigcup\{A\in\mathcal F:A\subseteq B\},\qquad
 e_{\mathcal F}(B)=\{j:j\in\bas_{\mathcal F}(B\cup\{j\})\}.
\]
The basis is the largest feasible subset of $B$. The set $e_{\mathcal F}(B)$
consists of that basis together with its available one-element
extensions. To verify this description, take a feasible subset of
$B\cup\{j\}$ containing $j$ and an accessible order of it. Its prefix
before $j$ is feasible inside $B$; unioning with the basis makes $j$
an available extension. The converse is immediate.

\begin{theorem}[Exact antimatroid representation]\label{reset:thm:anti}
Reset term functions on a fixed input set $N$ correspond bijectively to
antimatroids with one feasible singleton on nonempty supports
$X\subseteq N$. The correspondence is
\[
 E_t=e_{\mathcal F_t},\qquad
 \mathcal F_t=\{A\subseteq X:A\subseteq E_t(A)\}.
\]
For reset itself the family operation is
\begin{equation}\label{reset:eq:family-product}
 \mathcal F\del\mathcal G=
 \{\varnothing\}\cup
 \{A\cup C:A\in\mathcal F,\ C\in\mathcal G\setminus\{\varnothing\}\}.
\end{equation}
The root of the product is the root of $\mathcal G$.
\end{theorem}
\begin{proof}
The variable $x_j$ gives $\{\varnothing,\{j\}\}$.
The family in \eqref{reset:eq:family-product} is union-closed. To order a
nonempty member $A\cup C$ accessibly, first take an accessible order of
$C$ in $\mathcal G$, then one of $A$ in $\mathcal F$, skipping positions
already included. Every prefix is a union of a feasible set of
$\mathcal F$ with a nonempty feasible set of $\mathcal G$. This proves
accessibility, full support, and the unique initial element.

If $r$ is the root of $\mathcal G$, its basis obeys
\[
 \bas_{\mathcal F\del\mathcal G}(B)=
 \begin{cases}
 \varnothing,&r\notin B,\\
 \bas_{\mathcal F}(B)\cup\bas_{\mathcal G}(B),&r\in B.
 \end{cases}
\]
Applying the definition of $e$ gives exactly the selector recursion
in Lemma~\ref{reset:lem:selection}. Also
$A\subseteq e_{\mathcal F}(A)$ iff $A$ is feasible: membership of
$j\in A$ in this readiness set means $j\in\bas_{\mathcal F}(A)$.
This proves faithful representation for every term.

Conversely, an accessible order $r,j_2,\ldots,j_m$ of the full support
is a basic word. Its path term
\[
 x_{j_m}\del x_{j_{m-1}}\del\cdots\del x_r
\]
with \emph{left} parentheses has the feasible family of its prefixes.
Here the root $r$ is the last argument of the term, while it is the
first element of the accessible order. Combining two terms with the
same root by reset is pointwise union. On families it is union closure
of the two families: each already contains the common singleton root.
Every feasible set extends to a basic word, by accessibility and union
closure. Thus joining the path terms for all basic words gives exactly
the prescribed antimatroid. This is an application of the standard
basic-word representation, with reset supplying its composition.
\end{proof}

The parentheses in the path display are essential: for an order
$r,a,b$, the term is $(x_b\del x_a)\del x_r$. It selects the initial
segment through the first resolved input. Repeated labels in more
general terms are handled by \eqref{reset:eq:family-product}; replacing the
occurrence tree by a graph on variable names can incorrectly combine
paths from different occurrences.

\subsection{Polynomial recognition and synthesis}

The family representation can be exponentially large. Prerequisite
antichains retain a polynomial-size description for a given term.
They also have an exact realizability criterion.

Given a candidate root $r$ and minimal antichains $\mathcal G_j$,
call a set $A$ \emph{constructible} if its positions can be ordered
$a_1,\ldots,a_m$ such that for every $\ell$ some member of
$\mathcal G_{a_\ell}$ lies in $\{a_1,\ldots,a_{\ell-1}\}$.
This is checked greedily: start empty and add any enabled member of
$A$. Enabling is monotone, so a greedy choice cannot prevent success.

\begin{theorem}[Exact realizability]\label{reset:thm:synthesis}
A table is the canonical table of a reset term exactly when:
\begin{enumerate}
\item $\mathcal G_r=\{\varnothing\}$ and no other antichain contains
$\varnothing$;
\item every member of $\mathcal G_j$ omits $j$;
\item every member of every antichain is constructible.
\end{enumerate}
Recognition is polynomial in the table description. A realizable
table with $M$ prerequisites on $n$ input positions has a term
realization of size $O(nM)$.
\end{theorem}
\begin{proof}
For necessity, roots and self-independence give the first two
conditions. A minimal prerequisite $A$ enabling $j$ is feasible in the
represented antimatroid: readiness at $A$ has a feasible predecessor
inside $A$, and minimality makes that predecessor all of $A$.
Any accessible order of $A$ verifies constructibility.

For sufficiency, for each prerequisite $A\in\mathcal G_j$ choose its
greedy order $a_1,\ldots,a_m$ and make the path term with order
$a_1,\ldots,a_m,j$. Its last position is not repeated, by condition 2.
The first is the unique unconditionally enabled position $r$.
The endpoint has exactly prerequisite $A$. Every earlier position on
the path has a prerequisite contained in the preceding prefix, by
constructibility. Thus all the path's extra contributions are already
permitted by the candidate table. Join these paths by reset with their
common root. Every required prerequisite is present and no unwanted
contribution is introduced. There are $M$ paths of length at most
$n+1$.
\end{proof}

The constructibility condition has a simpler equivalent local form:
\begin{equation}\label{reset:eq:complete}
 A\in\mathcal G_j,\ k\in A
 \quad\Longrightarrow\quad
 \exists B\in\mathcal G_k\; B\subseteq A.
\end{equation}
Constructibility implies this completeness condition. Conversely,
induct on the size of a prerequisite $A$. For each $k\in A$, choose
$B\in\mathcal G_k$ inside $A$. Since $k\notin B$, it is smaller than
$A$ and constructible by induction. Append $k$ to its construction.
Unions of constructible sets are constructible: concatenate their
orders, skipping repeated positions. Unioning these constructions
reaches all of $A$.

On actual support, the antichain, nonemptiness, self-exclusion, and
completeness conditions are precisely the axioms for a \emph{dependency
structure with choice} (DSC) in
\cite[Definition~2.9]{BazermanMinichielloPuzio2023}. Their correspondence
with antimatroids is established in that work
\cite[Theorem~3.15]{BazermanMinichielloPuzio2023}. Thus the prerequisite
model reuses DSCs, with one unconditionally enabled event, as a concrete
presentation of reset term functions. The implementation checks the
local condition and uses greedy orders for synthesis.

An exponential antimatroid is therefore not needed for either
equality or synthesis. These are algorithms for formulas over the
primitive reset operation. Its expansion into owned choices may
duplicate subterms; no polynomial bound on that separate expansion
is claimed.

\subsection{Substitution, counting, and minimality}

Let an outer table have root $r$ and prerequisites $\mathcal G_i$.
Substitute inner tables with roots $r_i$ and prerequisites
$\mathcal H^i_j$. Since inner ambiguity is determined exactly by its
root, substitution gives
\begin{equation}\label{reset:eq:substitution}
 r_{\rm new}=r_r,\qquad
 \mathcal G^{\rm new}_j=
 \min\bigcup_i
 \left\{\{r_k:k\in A\}\cup C:
       A\in\mathcal G_i,\ C\in\mathcal H^i_j\right\}.
\end{equation}
Here $\min$ deletes nonminimal sets. The formula follows by substituting
the inner output unions into the outer output union, and substituting
each inner root status into an outer guard. It is polynomial in the
two explicit table descriptions. It respects actual term substitution,
including mixed inputs, rather than only a homogeneous observation.

Removing the unique initial element from every nonempty feasible set
leaves an arbitrary antimatroid on the remaining support. Conversely,
adjoining one initial element to every feasible set gives the unique-root
form. If $a_m$ counts labeled antimatroids on a full $m$-element support,
with $a_0=1$, the number of $n$-ary reset term functions is therefore
\[
 f_n=\sum_{m=1}^n\binom nm\,m\,a_{m-1}.
\]
Independent finite enumerations give
\[
 (a_0,\ldots,a_4)=(1,1,3,22,485),\qquad
 (f_1,\ldots,f_5)=(1,4,18,140,2980).
\]
The antimatroid counts are counts of known combinatorial objects;
the equation describes their exact relationship to reset operations.

\begin{theorem}[Identification-minor minimality]\label{reset:thm:minimal}
Every nonprojection reset term has a binary identification minor equal
to reset. Hence its term clone is minimal.
\end{theorem}
\begin{proof}
A term selects all its occurring variables when every input is
ambiguous. Each of them is essential, as witnessed by assigning a proper
ambiguous set to all inputs except one assigned a larger set.
A nonprojection therefore has support of size at least two. Identify
its root variable with $y$ and every other occurring variable with $x$.
The resulting binary term has root $y$ and both arguments essential.
Its $x$ prerequisite cannot contain $x$, cannot be empty, and can only
use $y$. Thus it is $\{\{y\}\}$, while the $y$ prerequisite is empty.
This is exactly $x\del y$. Minimality follows.
\end{proof}

Minimality is already known for $G_3$ and its dual in
Waldhauser's work. The direct proof here records the stronger explicit
recovery by input identification, using the prerequisite representation.
